\subsection{Implicit Domain of a Function} 
When we define a function using \makeblank[1.5in] or a \makeblank[1.5in], it is easy to see the domain. It is more difficult when we have a relation defined by an equation.
\begin{definition}[Implicit Domain]
The \term{implicit domain} of a function $f(x)=\mbox{[expression]}$ is the set of all real numbers for which $\mbox{[expression]}$ is defined.
\end{definition}
\begin{example}
Let $f(x)=\dfrac{1}{x}$. What is the implicit domain of $f$?\lbr
To evaluate $f(x)$ we must \makeblank[1in] by $x$. We can \makeblank[1in] by any number except for \makeblanksmall, so the implicit domain of $f$ is \makeblank.
\end{example}
\begin{note}Learning to find the implicit domains of various kinds of functions will be one of our major tasks throughout this course!\end{note}
We can also define an \term{explicit domain} for a function.
\begin{definition}[Explicit Domain]
The \term{explicit domain} of a function is a part of the implicit domain, explicitly written as part of the definition of the function. We write
$$f(x)=\mbox{[expression]},\ \mbox{[restriction]}$$
to indicate that the domain of the function $f$ is restricted.
\end{definition}
\begin{example}
Let $f(x)=x^2,\ x\geq 0$.
\begin{itemize}
\item \makeblankfill{$\dom f$ is }
\item \makeblankfill{$f(4)=$}
\item \makeblankfill{$f(-2)$ is \makeblank[1in] because }
\end{itemize}
\end{example}
If an explicit domain is not given, $\dom f$ means the implicit domain of $f$.